
\documentclass[a4paper,12pt]{article} 

\usepackage{graphicx}
\usepackage{enumitem}
\usepackage{tabulary}
\usepackage{color}
\usepackage{textcomp}
\usepackage{listings}
\begin{document}
\title{Practicing for a Coding Interview}
\maketitle
\newpage

\section*{Practicing for a Coding Interview}

\emph{Captured: 2016-11-01 at 19:03 from <a href="https://medium.com/code-school/practicing-for-a-coding-interview-bdbb983e10e7?source=userActivityShare-c79006fee040-1478023380">medium.com</a>}

\includegraphics[width=\textwidth]{presentationImg_1.jpeg}

Interviews for programming jobs aren’t like most interviews you hear about.
There’s less, “What is your greatest weakness?” and, “Why are you right person
for the job?” and way more code — \emph{lots} of code.

That’s right, software engineering interviews involve writing code on the
spot, right in front of your interviewer. And if that sounds terrifying to
you, you’re not alone.

The best way to conquer the fear of coding interviews is to get \emph{really,
really good} at them. And it can be done! It just takes careful practice. So
today I want to share 5 ways you can practice and some advice for acing each
and every interview.

\subsubsection*{1. Write code on a piece of paper or whiteboard.}

From my experience, most companies will have you write code on a whiteboard
during the onsite interview.

If you’re not used to coding on a whiteboard, it’s quite awkward at first. You
have to write out every single character, you can’t easily insert new lines or
delete whole blocks, and you don’t have the helpful highlighting you’re used
to from your favorite code editor.

So iron out that awkwardness \emph{before} your interviews. Run a few practice
problems on a piece of paper or, if you can, on a real whiteboard.

Once you have the code hand-written, try debugging it. Run through what it
does, step by step, and look for errors. Then, type it up and actually test
it.

\subsubsection*{2. Read up on language-specific trivia.}

Interviewers love asking trivia about coding languages. But don’t stress out —
they usually won’t throw you questions about languages you don’t know. The
general rule is that any language listed on your resume is fair game, but most
interviewers will focus on the 1 or 2 languages you state as your most
comfortable.

For web engineering jobs, folks \emph{love} to ask about JavaScript trivia. Last
time I interviewed, I was asked about <a href="https://developer.mozilla.org/en-US/docs/Web/JavaScript/Closures">JavaScript
closures</a> 3
times! People also like to ask about
<a href="http://www.w3schools.com/js/js_hoisting.asp">hoisting</a> and the <a href="https://developer.mozilla.org/en-
US/docs/Web/JavaScript/Guide/Details_of_the_Object_Model">JavaScript
object model</a>.

For Python folks, interviewers love to ask about the <a href="http://programmers.stackexchange.com/a/186905">global interpreter
lock</a>, or the<a href="http://sebastianraschka.com/Articles/2014_python_2_3_key_dif\kern0pt{}f.html">dif\kern0pt{}ference
between Python 2 and Python
3</a>. For
Java, it’s the <a href="https://www.interviewcake.com/java-interview-questions#strings">dif\kern0pt{}ference between string literals and String
objects</a>. And
for Ruby, it’s <a href="http://blog.jayfields.com/2008/03/ruby-inject.html">some fanciness with
.inject()</a>.

The idea behind the questions is the same — to see whether or not you
understand the quirks and specifics of the languages you use. This helps the
interviewer gauge if you’re the kind of engineer who thinks not just about
writing code, but about writing_good_ code.

The good news is there are plenty of language-specific coding interview
question lists available for free online, and a quick Google search will turn
up lots of options to look through and help prepare you.

\subsubsection*{3. Don’t be afraid of algorithmic thinking.}

Many interviewers like to focus their questions on algorithm design. What does
that mean exactly? That means writing code that not only works, but works
\emph{ef\kern0pt{}ficiently}. An important (and great) tool for this is <a href="https://en.wikipedia.org/wiki/Big_O_notation">Big O
notation</a> — and if it’s new to
you, you’ll want to <a href="https://www.interviewcake.com/article/big-o
-notation-time-and-space-complexity">practice it</a>.

It’s tempting to skip this algorithmic problem-solving piece altogether —
especially if it’s new to you — and I see so many candidates make this
mistake. And remember, it doesn’t matter if you’re interviewing for Google or
a small start-up — many small shops like to talk about algorithm design just
as much as the big guys.

If this stuf\kern0pt{}f is hard for you, you’re not alone. This is the hardest part of
coding interviews for \emph{most} candidates, even folks who have a computer
science degree. But you’ll find that Big O notation and algorithm design come
together more quickly than you’d think. On top of all that, and maybe most
importantly, it really will make you a better engineer. It totally changes
your relationship with coding. Eventually you’ll find that you can rigorously
reason through whether your code is the best solution to a problem.

\subsubsection*{4. Write down your mistakes.}

Grab a fresh notebook. After each question you try, look back and ask
yourself, “What did I get wrong about this problem at first?” In other words,
what did you learn? For each question, take the time to write down 1 or 2
things you learned.

After each practice session, read through your whole running list of
learnings. This adds a nice layer of rigor to your practice, so you remember
things to fix for next time. You’ll be less likely to repeat the same mistakes
and, eventually, checking for those things will become second nature.

\subsubsection*{5. Schedule lower-stake interviews first!}

The very best practice for coding interviews is other coding interviews. The
more you do, the better you’ll be. So, if possible, don’t interview with your
#1 choice company right away — schedule a few others first to maximize the practice you get ahead of time.

Apply to several companies at once. It means more practice, and if you’re
lucky enough to get multiple of\kern0pt{}fers, it gives you more negotiating leverage!
Lots of companies can be found and contacted quickly on sites like
<a href="https://www.whitetruf\kern0pt{}fle.com/">WhiteTruf\kern0pt{}fle</a>,
<a href="https://angel.co/jobs">AngelList</a>, and <a href="https://hired.com/">Hired</a>. But
remember, don’t overdo it — you don’t want to waste anyone’s time. Let
yourself be open to opportunities you might not initially consider during your
search (after all, my last job was at a company I wouldn’t have considered if
I hadn’t found them on AngelList!). Most importantly, just remember to have
confidence in yourself and your skills — pair that with plenty of practice,
and interviewing will be a piece of cake.
\end{document}
